\documentclass[11pt]{article}
\usepackage[T1]{fontenc}
\usepackage{lmodern,microtype}
\usepackage[margin=1in]{geometry}
\usepackage{amsmath,amssymb,amsthm,mathtools,mathrsfs}
\usepackage{booktabs,tabularx,array}
\usepackage[dvipsnames]{xcolor}
\usepackage[colorlinks=true,linkcolor=MidnightBlue,citecolor=MidnightBlue,urlcolor=MidnightBlue]{hyperref}
\usepackage{bookmark,enumitem}
\setlist{nosep}
\setcounter{tocdepth}{1}
\setlength{\emergencystretch}{2em}
\renewcommand{\arraystretch}{1.15}
\newtheorem{theorem}{Theorem}[section]
\newtheorem{proposition}[theorem]{Proposition}
\newtheorem{lemma}[theorem]{Lemma}
\theoremstyle{definition}
\newtheorem{definition}[theorem]{Definition}
\theoremstyle{remark}
\newtheorem{remark}[theorem]{Remark}
\numberwithin{equation}{section}
\DeclareMathOperator{\Tr}{Tr}
\DeclareMathOperator{\Dom}{Dom}
\DeclareMathOperator{\HS}{HS}
\DeclareMathOperator{\Scal}{Scal}
\newcommand{\R}{\mathbb R}
\newcommand{\C}{\mathbb C}
\newcommand{\K}{\mathcal K}
\newcommand{\Href}{\mathcal K_{\mathrm{ref}}}
\newcommand{\HH}{\mathscr H}
\newcommand{\Dc}{\mathcal D_C}
\newcommand{\Obs}{\mathcal O}
\newcommand{\dd}{\,\mathrm d}
\newcommand{\ee}{\mathrm e}
\newcommand{\ii}{\mathrm i}
\newcommand{\norm}[1]{\left\lVert#1\right\rVert}
\newcommand{\abs}[1]{\left|#1\right|}
\newcommand{\braket}[2]{\left\langle#1,#2\right\rangle}
\newcolumntype{Y}{>{\raggedright\arraybackslash}X}
\title{Hamiltonian heat evolution of theta amplitudes\\[4pt]
\large Complete spectral realization and calibrated Newman response}
\author{Drafted for Edward Baker}
\date{6 November 2026\\\small Working manuscript}
\hypersetup{pdftitle={Hamiltonian heat evolution of theta amplitudes},pdfauthor={Drafted for Edward Baker},pdfsubject={Analytical results for Hamiltonian auxiliary dynamics}}
\begin{document}
\maketitle
\begin{abstract}
We develop an auxiliary Hamiltonian representation of the canonical
de Bruijn--Newman transform. The squared geometric Dirac operator on
$\R\times S^2_{1/2}$ has the complete dispersion $k^2+4n^2$; its radial
Green diagonal supplies an angular weight that cancels the spinor
multiplicities. A prescribed damping scan and differential processing
produce the canonical theta source. We use this source to prepare an
operator-valued amplitude, evolve it by the left-acting heat semigroup,
and pair it with a matched passive reference. The absolutely calibrated
response is
$\Obs_{\lambda,\tau}(z)=
 [Q_1(\pi+4\tau)/Q_1(\pi)]\ee^{-\tau z^2}H_\lambda(z)$.
We derive the self-adjoint realizations, exact prepared graph norms,
complete source and spectral tails, weighted heat estimates, and every
fixed finite mixed derivative, including actual right derivatives at
zero proper time. The nonvanishing multiplier preserves zeros and their
multiplicities and conjugates the Newman equation, while its inverse
amplifies high real frequencies without bound. These are analytical
statements about a supplied preparation and phase-sensitive readout;
they do not establish intrinsic preparation, a laboratory implementation,
or a collision exclusion. The final section identifies the pending
conditioning analysis and the source-specific estimates needed for
further progress.
\end{abstract}
\tableofcontents
\clearpage

\section{The construction and its analytical content}
\label{sec:overview}

The central object is the even entire function
\begin{equation}
 H_\lambda(z)=\int_0^\infty
 \ee^{\lambda v^2}\Phi(v)\cos(zv)\dd v,
 \qquad
 \Phi(v)=\sum_{m\ge1}
 \bigl(2\pi^2m^4\ee^{9v}-3\pi m^2\ee^{5v}\bigr)
 \ee^{-\pi m^2\ee^{4v}},\quad v\ge0.
 \label{eq:intro-H}
\end{equation}
The source has a smooth even extension to the full line. Its rapid decay
makes every real deformation parameter $\lambda$ admissible. With the
normalization in \eqref{eq:intro-H},
$H_0(z)=\xi((1+\ii z)/2)/8$ and
$\partial_\lambda H_\lambda=-\partial_z^2H_\lambda$
\cite{RodgersTao2020}. We will derive both identities from the source.

The auxiliary construction has three layers. First, the complete cylinder
Hamiltonian supplies a Green weight and a heat semigroup. Second, a
specified scan of the angular damping produces $\Phi$, which is placed
as a spatial amplitude. Third, an explicitly calibrated linear readout
observes the evolved amplitude. Keeping these layers distinct isolates
what the geometry proves from what the preparation and detector supply.

The resulting theorem is simple to state. Write
\begin{equation}
 Q_1(\beta)=\sum_{n\ge1}\ee^{-\beta n^2},\qquad
 q(\tau)=\frac{Q_1(\pi+4\tau)}{Q_1(\pi)}.
 \label{eq:intro-q}
\end{equation}
The normalized angular bank, the prepared field, and the readout are
defined precisely in Sections~\ref{sec:preparation} and~\ref{sec:readout}.

\begin{theorem}[Complete response, overview]
\label{thm:overview}
For the specified source preparation and matched linear readout, the
left-acting heat evolution on the complete cylinder gives
\begin{equation}
 \boxed{\Obs_{\lambda,\tau}(z)
       =q(\tau)\ee^{-\tau z^2}H_\lambda(z).}
 \label{eq:response-overview}
\end{equation}
This holds for every real $\lambda$, every $\tau\ge0$, and every $z\in\C$.
For any fixed finite $M,T,Z,b\ge0$, all fixed finite mixed derivatives
in $\lambda,\tau,z$ are continuous on
\begin{equation}
 |\lambda|\le M,\quad 0\le\tau\le T,\quad
 |\Re z|\le Z,\quad |\Im z|\le b,
 \label{eq:windows}
\end{equation}
and can be passed through the complete source, heat, Fourier, and trace
operations. Proper-time derivatives at zero are right derivatives.
Every bound depends on the specified finite derivative orders and window.
\end{theorem}

We first realize the complete self-adjoint operators, then derive
summable angular tails and integrable radial majorants, and finally
evaluate the response. This order also explains why fixed-frequency
observation, although not bounded on a generic radial $L^2$ space, is
well-defined on the prepared and evolved family.

There are several distinct parameters:
\begin{center}
\begin{tabularx}{\textwidth}{@{}lY@{}}
\toprule
Parameter & Meaning \\
\midrule
$v$ & Dimensionless physical cylinder coordinate. \\
$u$ & External scan coordinate, later imprinted as $u=v$. \\
$\beta$ & Angular damping parameter; the source scan uses $\beta=\pi\ee^{4u}$. \\
$\lambda$ & Newman parameter in the prepared source $\ee^{\lambda v^2}\Phi(v)$. \\
$\tau$ & Nonnegative dimensionless heat or Euclidean proper time. \\
$k$, $z$ & Real radial momentum and its complex continuation. \\
\bottomrule
\end{tabularx}
\end{center}
In particular, increasing $\tau$ is not increasing $\lambda$. Ordinary
unitary laboratory evolution is a third evolution, discussed below.

The parent manuscript provides the common background for the direct and
auxiliary branches \cite{TNDBackground}. The present manuscript develops
the auxiliary branch on its own. Its principal mathematical input is the
complete-spectrum response note \cite{TNDAux}, accepted within its
conditional scope by the branch assessment \cite{TNDAuxReview}, together
with the shared geometry and source analysis
\cite{TNDParent,TNDSharedReview}. The first-jet conditioning program
\cite{TNDAuxConditioning} remains pending; it is separated from the
established results in Section~\ref{sec:outstanding}.

\section{Cylinder geometry, complete spectrum, and the Green weight}
\label{sec:geometry}

\subsection{The geometric square}

Let $\K=L^2(S^2,\Sigma;\dd\Omega)$ be the unit-sphere spinor space.
Use the self-adjoint convention $D=-\ii\not{\nabla}$ and let
$\Gamma$ be parallel chirality. Then
\begin{equation}
 \Gamma^*=\Gamma,\qquad \Gamma^2=I,\qquad
 \Gamma D=-D\Gamma,
 \qquad D\chi_{n,a}^{\pm}=\pm n\chi_{n,a}^{\pm},
 \quad 1\le a\le2n,\quad n\ge1.
 \label{eq:sphere-spectrum}
\end{equation}
We take $\chi^-_{n,a}=\Gamma\chi^+_{n,a}$ and write $A=D^2$.
The projection $P_n$ onto the complete $A=n^2$ shell has rank $4n$.
These eigenvalues and multiplicities are the $S^2$ specialization of the
spinor spectrum in \cite[eqs.~(3.40), (3.54)]{CamporesiHiguchi1996}.

Consider the complete cylinder
\begin{equation}
 C=\R\times S^2_{1/2},\qquad
 g_C=\dd v^2+\tfrac14g_{S^2},\qquad
 \dd\mathrm{vol}_C=\tfrac14\dd v\dd\Omega.
\end{equation}
In the product-measure representation its geometric Dirac operator is
\begin{equation}
 \Dc=-\ii\Gamma\partial_v+2D,
 \qquad L_0=\Dc^2=-\partial_v^2+4A.
 \label{eq:cylinder-square}
\end{equation}
The cross terms cancel by the anticommutation in
\eqref{eq:sphere-spectrum}. Multiplication by $2$ is the unitary map from
product-measure spinors to metric-volume spinors. It does not change the
differential expression.

The cylinder scalar curvature is $8$. The Lichnerowicz identity therefore
also reads $\Dc^2=\nabla^*\nabla+2$. This is the same operator as
\eqref{eq:cylinder-square}: the curvature contribution is already inside
the Dirac square, and adding a further $2$ would change the model.
The slash and curvature conventions agree with
\cite[eqs.~(5.35)--(5.36)]{CamporesiHiguchi1996} after multiplication by
$-\ii$ to make the Dirac operator self-adjoint.

For a length scale $\ell_0>0$ and mass $m_{\mathrm{phys}}>0$, the supplied
energy convention is
\begin{equation}
 H_{\mathrm{kin}}=\varepsilon L_0,
 \qquad \varepsilon=\frac{\hbar^2}{2m_{\mathrm{phys}}\ell_0^2},
 \qquad \tau=\frac{\varepsilon t_E}{\hbar}.
 \label{eq:physical-scale}
\end{equation}
Thus $\ee^{-\tau L_0}$ is Euclidean evolution. Real laboratory time would
instead give $\ee^{-\ii\varepsilon t_{\mathrm{lab}}L_0/\hbar}$.

\subsection{Self-adjoint realization without a radial cutoff}

Use the unitary full-line Fourier convention
\begin{equation}
 \widehat f(k)=(2\pi)^{-1/2}\int_\R\ee^{-\ii kv}f(v)\dd v.
 \label{eq:fourier}
\end{equation}
In each pair $(\chi^+_{n,a},\chi^-_{n,a})$ the first-order multiplier is
\begin{equation}
 M_n(k)=\begin{pmatrix}2n&k\\k&-2n\end{pmatrix},\qquad
 E_n(k)=\sqrt{k^2+4n^2},\qquad
 \Pi_\pm(n,k)=\tfrac12\left(I\pm\frac{M_n(k)}{E_n(k)}\right).
 \label{eq:signed-block}
\end{equation}
Since $M_n(k)^2=\omega_n(k)I$ with $\omega_n(k)=k^2+4n^2$,
the maximal domains on $L^2(\R;\K)$ are
\begin{align}
 \Dom\Dc=\Dom L_0^{1/2}
 &=\left\{f:\sum_n\int_\R\omega_n(k)
           \norm{P_n\widehat f(k)}^2\dd k<\infty\right\},
 \label{eq:dirac-domain}\\
 \Dom L_0^p
 &=\left\{f:\sum_n\int_\R\omega_n(k)^{2p}
           \norm{P_n\widehat f(k)}^2\dd k<\infty\right\},
 \qquad p\ge0.
 \label{eq:square-domain}
\end{align}

\begin{proposition}[Complete realization]
\label{prop:complete-realization}
The multipliers above define self-adjoint operators. Finite complete-shell
sums with radial coefficients in $C_c^\infty(\R)$ are cores for $\Dc$
and every fixed finite integer power of $L_0$. Their spectra are
\begin{equation}
 \sigma(\Dc)=(-\infty,-2]\cup[2,\infty),\qquad
 \sigma(L_0)=[4,\infty),
 \label{eq:spectra}
\end{equation}
and are purely absolutely continuous in the radial spectral variable.
\end{proposition}

\begin{proof}
Hermitian matrix multiplication is symmetric on its maximal domain.
Testing a vector in the adjoint domain against compactly supported
spectral functions in finitely many channels forces its image under
$M_n(k)$ to be square-integrable. This gives exactly
\eqref{eq:dirac-domain}; the scalar argument gives
\eqref{eq:square-domain}. Hence the maximal operators equal their
adjoints.

Projection onto $\sum_{n\le N}P_n$ converges in each displayed graph norm
by dominated convergence and commutes with the signed blocks. On this
finite angular space the first-order graph norm is equivalent to the
radial $H^1$ norm, and the graph norm of $L_0^p$ to $H^{2p}$ for integer
$p$. Smooth cutoffs and mollification give the stated cores. No radial
boundary condition is introduced.

The eigenvalues of \eqref{eq:signed-block} are $\pm E_n(k)$, each of
physical multiplicity $2n$; the square has multiplicity $4n$. Their
essential ranges give \eqref{eq:spectra}. On each open half-line in $k$,
the dispersions are smooth and strictly monotone with nonzero derivative.
Pushing an integrable spectral density through them gives an absolutely
continuous measure. The exceptional threshold preimages are singletons
and have measure zero. Countable sums preserve absolute continuity.
\end{proof}

In particular, a fixed radial plane wave, cosine, or threshold value is
not an $L^2$ eigenstate. Restricting to even radial functions reduces the
square, but not the first-order operator: its radial derivative changes
parity. Complete shells are the natural cutoffs when both signs of the
Dirac operator are retained.

\subsection{The Green diagonal and multiplicity cancellation}

On shell $n$, the inverse of $-\partial_v^2+4n^2$ has kernel
\begin{equation}
 g_n(v,w)=\frac{\ee^{-2n|v-w|}}{4n}.
 \label{eq:green-kernel}
\end{equation}
Indeed, away from $v=w$ it solves the homogeneous equation, while its
first derivative has jump $-1$, producing the required unit delta
under $-\partial_v^2$. Equivalently,
\begin{equation}
 g_n(v,v)=\int_\R\frac{\dd k}{2\pi(k^2+4n^2)}=\frac1{4n}.
 \label{eq:green-diagonal}
\end{equation}
This is the complete radial integral. Evaluation of the resolvent at
$k=0$ would give $1/(4n^2)$ and would not yield this weight.

The angular operator supplied by the diagonal is
\begin{equation}
 W=(4\sqrt A)^{-1}=\sum_{n\ge1}\frac1{4n}P_n.
 \label{eq:green-weight}
\end{equation}
It is bounded and positive but not trace class. For every $\beta>0$,
however,
\begin{equation}
 \Tr_\K(W\ee^{-\beta A})
 =\sum_{n\ge1}(4n)\frac1{4n}\ee^{-\beta n^2}
 =Q_1(\beta).
 \label{eq:theta-trace}
\end{equation}
This exact cancellation is the link between the cylinder geometry and
the unit coefficients of the theta series. It retains all shells and
both radial directions. It is not the trace of the full cylinder heat
operator, which has infinite radial volume.

\section{From the angular trace to the canonical source}
\label{sec:source}

\subsection{The scan, differential processing, and reflection}

The trace \eqref{eq:theta-trace} becomes a spatial profile only after an
external prescription. Set
\begin{equation}
 \beta(u)=\pi\ee^{4u},\qquad
 p(u)=\ee^uQ_1(\pi\ee^{4u}),\qquad
 \Phi(u)=\tfrac18(\partial_u^2-1)p(u).
 \label{eq:scan}
\end{equation}
For fixed $u$, differentiation of the Gaussian series is absolute and
locally uniform. With $x=\pi\ee^{4u}$ it gives
\begin{equation}
 \Phi(u)=\ee^u\sum_{m\ge1}
 \bigl(2(m^2x)^2-3m^2x\bigr)\ee^{-m^2x}.
 \label{eq:phi-series}
\end{equation}
The source summation index is written $m$ because the scalar source,
once prepared, is independent of the physical channel index $n$ of
the evolving field. Equating the scan coordinate $u$ to the physical
coordinate $v$ is also part of the preparation. The translation-invariant
Green kernel alone does not prescribe either operation.

The complete theta identity is what makes this source smooth and even.
For $R>0$, let
\begin{equation}
 \Theta(R)=\sum_{j\in\mathbb Z}\ee^{-\pi j^2R^2}
          =1+2Q_1(\pi R^2).
\end{equation}
The modular identity
$\Theta(R)=R^{-1}\Theta(R^{-1})$ is the zero-argument specialization
of \cite[20.7.32]{DLMFThetaModular}. Taking $R=\ee^{2u}$ yields
\begin{equation}
 p(u)-p(-u)=-\sinh u.
 \label{eq:reflection-p}
\end{equation}
The operator $\partial_u^2-1$ annihilates the right side, so
\begin{equation}
 \Phi(u)=\Phi(-u).
 \label{eq:phi-even}
\end{equation}
The series defining $p$ is smooth on every compact interval of the full
line; thus this is a smooth reflection, including the origin. Both the
zero-mode correction in $\Theta$ and the complete unit coefficients are
essential to this argument.

Every summand of \eqref{eq:phi-series} is positive when $u\ge0$, since
$m^2x\ge\pi>3/2$. Therefore $\Phi>0$ on the whole line. The first term
also gives, with $x=\pi\ee^{4|v|}$,
\begin{equation}
 \Phi(v)\ge\left(2-\frac3\pi\right)
 \ee^{|v|}x^2\ee^{-x}>0.
 \label{eq:phi-lower}
\end{equation}
This positivity concerns the source. It will not determine the sign of
its oscillatory Fourier observation.

\subsection{A hierarchy of complete-source bounds}

We next establish a single collection of majorants that will control the
preparation, evolution, and readout. Define polynomials
\begin{equation}
 P_0(x)=2x^2-3x,\qquad
 P_{j+1}(x)=(1-4x)P_j(x)+4xP_j'(x).
 \label{eq:source-poly}
\end{equation}
Then $\deg P_j=j+2$, and differentiating on $u\ge0$ gives
\begin{equation}
 \Phi^{(j)}(u)=\ee^u\sum_{m\ge1}
                 P_j(m^2x)\ee^{-m^2x}.
 \label{eq:phi-derivatives}
\end{equation}
Let $a_j$ be the sum of the absolute coefficients of $P_j$, and put
\begin{equation}
 C_d(N)=\sum_{m>N}m^{2d}\ee^{-\pi(m^2-1)},\qquad
 B_j=a_jC_{j+2}(0).
 \label{eq:source-constants}
\end{equation}
Every $C_d(N)$ is finite and tends to zero with $N$. Since $x\ge\pi>1$,
\begin{equation}
 \ee^u\sum_{m>N}|P_j(m^2x)|\ee^{-m^2x}
 \le a_j C_{j+2}(N)\ee^u x^{j+2}\ee^{-x}.
 \label{eq:source-tail}
\end{equation}
Indeed, $|P_j(m^2x)|\le a_jm^{2(j+2)}x^{j+2}$ and
$\ee^{-(m^2-1)x}\le\ee^{-\pi(m^2-1)}$. Thus the estimate controls
the sum of absolute values, not merely cancellation in a signed tail.
Reflection gives the global bound
\begin{equation}
 |\Phi^{(j)}(v)|\le B_j\ee^{|v|}x^{j+2}\ee^{-x},
 \qquad x=\pi\ee^{4|v|}.
 \label{eq:phi-global}
\end{equation}

For nonnegative integers $r,d$ and finite $M,b\ge0$, define
\begin{equation}
 J_{r,d}(M,b)=\int_\R |v|^r
 \ee^{Mv^2+(b+1)|v|}x^d\ee^{-x}\dd v.
 \label{eq:J-definition}
\end{equation}
These integrals are finite. More explicitly, the substitution
$x=\pi\ee^{4|v|}$ gives
\begin{equation}
 J_{r,d}(M,b)=\frac{2}{4^{r+1}}\int_\pi^\infty
 [\log(x/\pi)]^r(x/\pi)^{(b+1)/4}x^{d-1}
 \exp\!\left(\frac{M}{16}\log^2(x/\pi)-x\right)\dd x.
 \label{eq:J-integral}
\end{equation}
The term $-x$ dominates the logarithmic square and all polynomial factors.
The remaining integral beyond a radial cutoff consequently tends to
zero. This provides a complete radial tail bound at every fixed order.

Write
\begin{equation}
 F_\lambda(v)=\ee^{\lambda v^2}\Phi(v),\qquad
 f_{a,\lambda}(v)=v^{2a}F_\lambda(v)=\partial_\lambda^aF_\lambda(v).
 \label{eq:deformed-source}
\end{equation}
Define $R_0(\lambda,v)=1$ and
$R_{j+1}=\partial_vR_j+2\lambda vR_j$, so that
$\partial_v^j\ee^{\lambda v^2}=\ee^{\lambda v^2}R_j(\lambda,v)$.
For each fixed $a,h$ and $M$, finite Leibniz expansion gives
\begin{equation}
 |f_{a,\lambda}^{(h)}(v)|
 \le C_{a,h,M}(1+|v|)^{2a+h}
 \ee^{Mv^2+|v|}x^{h+2}\ee^{-x},\qquad |\lambda|\le M.
 \label{eq:f-majorant}
\end{equation}
For specificity, let $r_j(M)$ be the sum of the absolute coefficients
of $R_j$ after replacing every power $\lambda^t$ by $M^t$. One permissible
constant is
\begin{equation}
 C_{a,h,M}=\sum_{\substack{i+j+\ell=h\\i\le2a}}
 \frac{h!}{i!j!\ell!}\frac{(2a)!}{(2a-i)!}\,r_j(M)B_\ell.
 \label{eq:f-constant}
\end{equation}
Here $|R_j(\lambda,v)|\le r_j(M)(1+|v|)^j$, and
$x^{\ell+2}\le x^{h+2}$; these inequalities prove
\eqref{eq:f-majorant} directly.

It is useful to collect the weighted $L^1$ bounds as
\begin{align}
 w_{r,b}(v)&=(1+|v|)^r\ee^{b|v|},\\
 \mathcal N_{a,h;r,b}(M)
 &=\sup_{|\lambda|\le M}\norm{w_{r,b}f_{a,\lambda}^{(h)}}_1
 \nonumber\\
 &\le C_{a,h,M}\sum_{\ell=0}^{2a+h+r}
       \binom{2a+h+r}{\ell}J_{\ell,h+2}(M,b)<\infty.
 \label{eq:N-bound}
\end{align}
Squaring \eqref{eq:f-majorant} proves all finite $L^2$ derivative
bounds. Multiplying it by any polynomial and taking a supremum proves
every fixed Schwartz seminorm. The same majorants, with one further
factor $v^2$, prove continuity and differentiation in $\lambda$ in
each of these spaces. In particular, the image of a compact $\lambda$
interval is compact in every one of the specified norms.

Two distinctions about completeness will be used later. First, a
reflected finite source prefix $\Phi_N$ obeys the value-and-moment bound
\begin{equation}
 \norm{w_{r,b}v^{2a}\ee^{\lambda v^2}(\Phi-\Phi_N)}_1
 \le a_0C_2(N)\sum_{\ell=0}^r\binom r\ell
                      J_{2a+\ell,2}(M,b)\longrightarrow0.
 \label{eq:prefix-value-tail}
\end{equation}
Such a prefix generally has derivative jumps at the reflection point.
It is not a global smooth graph-core approximation; the operator cores
of Proposition~\ref{prop:complete-realization} supply those approximations.
Second, integration by parts yields, for any finite $d,l\ge0$,
\begin{equation}
 |\partial_k^d\widehat f_{a,\lambda}(k)|
 \le\frac{(1+k^2)^{-l}}{\sqrt{2\pi}}
 \norm{(1-\partial_v^2)^l[(-\ii v)^df_{a,\lambda}]}_1.
 \label{eq:fourier-tail}
\end{equation}
The right side is uniformly finite by \eqref{eq:N-bound}. Choosing $l$
sufficiently large controls every fixed polynomially weighted radial
Fourier tail in $L^1$ or $L^2$.

\subsection{The Newman transform and its absolute normalization}

Evenness converts the cosine transform into
\begin{equation}
 H_\lambda(z)=\frac12\int_\R F_\lambda(v)\ee^{-\ii zv}\dd v,
 \qquad \widehat F_\lambda(k)=\sqrt{\frac2\pi}\,H_\lambda(k).
 \label{eq:H-fourier}
\end{equation}
On $|\lambda|\le M$, $|\Im z|\le b$, the source estimates justify
every fixed finite mixed derivative:
\begin{equation}
 \partial_\lambda^a\partial_z^d H_\lambda(z)
 =\frac12\int_\R v^{2a}(-\ii v)^dF_\lambda(v)
                   \ee^{-\ii zv}\dd v,
 \qquad
 |\partial_\lambda^a\partial_z^d H_\lambda(z)|
 \le\frac{B_0}{2}J_{2a+d,2}(M,b).
 \label{eq:H-derivatives}
\end{equation}
The same domination on compact complex $\lambda$ sets shows joint
entire dependence on $(\lambda,z)$. In particular,
\begin{equation}
 \partial_\lambda H_\lambda=-\partial_z^2H_\lambda.
 \label{eq:newman-heat}
\end{equation}
The minus sign comes from $(-\ii v)^2=-v^2$. The parameter $\lambda$
reweights the source; it is not the cylinder heat parameter.

For completeness, the factor $1/8$ in the zeta normalization can be
derived directly, rather than left as an adjustable gain. By evenness we
may use $\ee^{\ii zv}$ in \eqref{eq:H-fourier}. Put
$s=(1+\ii z)/2$ and $t=\ee^{4v}$. For $\Re s>1$, absolute summability
allows termwise integration over the full line and gives
\begin{align}
 H_0(z)
 &=\frac18\sum_{m\ge1}\int_0^\infty
 t^{s/2-1}\bigl[2(\pi m^2t)^2-3\pi m^2t\bigr]
                         \ee^{-\pi m^2t}\dd t
 \nonumber\\
 &=\frac{2\Gamma(s/2+2)-3\Gamma(s/2+1)}8
                         \pi^{-s/2}\zeta(s)
 \nonumber\\
 &=\frac{s(s-1)}{16}\Gamma(s/2)\pi^{-s/2}\zeta(s)
  =\frac18\xi(s).
 \label{eq:xi-normalization}
\end{align}
Absolute integrals of the individual terms are bounded by a constant
times $m^{-\Re s}$ in this half-plane. The already established entire
integral supplies the continuation of the identity to all $z$.
This agrees with \cite[Introduction, eqs.~(1)--(4)]{RodgersTao2020}.
For real $\lambda$, $H_\lambda(0)>0$ by source positivity, so none of
these entire functions is identically zero.

\section{A complete operator-valued preparation}
\label{sec:preparation}

\subsection{The passive bank and left amplification}

Let $\Href$ be a passive copy of $\K$, and fix a matched unitary
$J:\Href\to\K$. With $Q=Q_1(\pi)$, define
\begin{equation}
 R=\frac{W\ee^{-\pi A}}Q=\sum_{n\ge1}r_nP_n,
 \qquad r_n=\frac{\ee^{-\pi n^2}}{4nQ},
 \qquad B=R^{1/2}J.
 \label{eq:bank}
\end{equation}
The same multiplicity cancellation as in \eqref{eq:theta-trace} gives
\begin{equation}
 \Tr R=1,\qquad BB^*=R,\qquad B^*B=J^*RJ,
 \qquad \norm B_{\HS}=1.
 \label{eq:bank-normalization}
\end{equation}
Every $r_n$ is positive, although they tend to zero, so $R$ is faithful
without a uniform positive lower bound. This is an angular normalization
only. It does not normalize the scalar amplitude $F_\lambda$.

For each finite $s\ge0$, put
\begin{equation}
 \mu_s=\frac1Q\sum_{n\ge1}n^{2s}\ee^{-\pi n^2}.
 \label{eq:mu}
\end{equation}
Then
\begin{align}
 \Tr(A^sR)&=\mu_s,\qquad
 \norm{A^sB}_{\HS}^2=\mu_{2s},
 \label{eq:bank-moments}\\
 \norm{A^s(I-P_{\le N})B}_{\HS}^2
 &=\frac1Q\sum_{n>N}n^{4s}\ee^{-\pi n^2}\longrightarrow0.
 \label{eq:bank-tail}
\end{align}
All these sums converge by Gaussian decay. The fixed reference is an
apparatus label, with no assigned right-hand energy.

The field space is
\begin{equation}
 \HH=L^2\bigl(\R;\HS(\Href,\K)\bigr)
 \simeq L^2(\R;\K)\otimes\overline{\Href}.
 \label{eq:amplified-space}
\end{equation}
Let $\Dc^{\mathrm L}$ and $L=L_0^{\mathrm L}$ act on the physical side,
amplified by the identity on the reference. The signed blocks
\eqref{eq:signed-block}, spectra \eqref{eq:spectra}, and domain formulas
\eqref{eq:dirac-domain}--\eqref{eq:square-domain} remain valid with
$\norm{P_n\widehat f}^2$ replaced by
$\norm{P_n\widehat S}_{\HS}^2$. Testing the multiplier adjoint in
finitely many reference columns proves self-adjointness exactly as before.
Complete physical shells, finite reference columns, and radial
$C_c^\infty$ functions form a common core. Reference projections converge
in graph norm because the squared column norms are summable. No extra
factor $4n$ or right-energy weight is introduced by the amplification.

\subsection{Exact preparation and graph norms}

The supplied initial amplitude is
\begin{equation}
 S_{\lambda,0}(v)=F_\lambda(v)B.
 \label{eq:prepared-field}
\end{equation}
It is a linear amplitude, rather than $\sqrt{F_\lambda}$ or a normalized
radial density. Its exact normalization is needed for
\eqref{eq:xi-normalization}. The spatial placement and the matched bank
in \eqref{eq:prepared-field} are model data.

\begin{proposition}[Prepared graph identities]
\label{prop:graph-identity}
For every finite integer $a,h,p\ge0$,
\begin{equation}
 \boxed{\norm{L^p(f_{a,\lambda}^{(h)}B)}_{\HH}^2
 =\sum_{j=0}^{2p}\binom{2p}{j}4^{2p-j}\mu_{2p-j}
                 \norm{f_{a,\lambda}^{(h+j)}}_2^2.}
 \label{eq:graph-identity}
\end{equation}
The expression is finite and uniform on compact real $\lambda$
intervals. The prepared family is differentiable to every finite order
in each of these graph norms.
\end{proposition}

\begin{proof}
By \eqref{eq:bank},
$\norm{P_nB}_{\HS}^2=\ee^{-\pi n^2}/Q$. The spectral domain norm is
therefore
\begin{equation*}
 \frac1Q\sum_n\ee^{-\pi n^2}\int_\R
 (k^2+4n^2)^{2p}k^{2h}|\widehat f_{a,\lambda}(k)|^2\dd k.
\end{equation*}
Expand the nonnegative polynomial by the finite binomial formula.
The angular sum in its $j$th term is $4^{2p-j}\mu_{2p-j}$, and
Plancherel turns the radial integral into
$\norm{f_{a,\lambda}^{(h+j)}}_2^2$. The bounds in
\eqref{eq:f-majorant} prove finiteness and continuity, including one
further parameter derivative whenever differentiation is required.
\end{proof}

This identity separates the two kinds of regularity: angular Gaussian
moments and radial source derivatives. It also provides the precise
prepared-data estimates at zero heat time, where generic smoothing
bounds would diverge. In conjunction with
\eqref{eq:bank-tail} and \eqref{eq:fourier-tail}, it controls both
physical angular and radial spectral tails. In a reference basis
matched by $J$, exhausting complete shells gives the same Gaussian
reference tails; in any fixed basis, convergence follows from the
summable column norms. The source index, physical shell, reference
column, and continuous momentum remain separate variables.

\section{Heat evolution with complete derivative control}
\label{sec:heat}

\subsection{Semigroup and prepared strong solutions}

Functional calculus for the realization in Section~\ref{sec:preparation}
gives
\begin{equation}
 P_n\widehat{\ee^{-\tau L}S}(k)
 =\ee^{-\tau(k^2+4n^2)}P_n\widehat S(k),\qquad
 \norm{\ee^{-\tau L}}=\ee^{-4\tau},\quad \tau\ge0.
 \label{eq:semigroup}
\end{equation}
Dominated convergence proves strong continuity at zero. For $r\ge0$ and
$\tau>0$,
\begin{equation}
 \norm{L^r\ee^{-\tau L}}
 =\sup_{x\ge4}x^r\ee^{-\tau x}<\infty.
 \label{eq:smoothing}
\end{equation}
This gives positive-time left-energy smoothing. It supplies neither
uniform arbitrary-data smoothing at zero nor reference-side regularity.
The semigroup is not compact: even one angular channel contains a
nonzero multiplication operator on a nonatomic radial spectral space.
In particular it is not trace class.

For data in $\Dom L$, the semigroup is the unique strong solution of
$\partial_\tau S=-LS$. Indeed, a difference $U$ of two strong solutions
with zero initial data satisfies
\begin{equation*}
 \frac\dd{\dd\tau}\norm{U(\tau)}^2
 =-2\braket{U(\tau)}{LU(\tau)}
 \le-8\norm{U(\tau)}^2,
\end{equation*}
so its norm remains zero. Define the prepared evolution by
\begin{equation}
 S_{\lambda,\tau}=\ee^{-\tau L}(F_\lambda B).
 \label{eq:evolved-field}
\end{equation}
Proposition~\ref{prop:graph-identity} gives, at every finite order,
\begin{align}
 \partial_\lambda^a\partial_v^h\partial_\tau^cS_{\lambda,\tau}
 &=(-1)^c\ee^{-\tau L}L^c(f_{a,\lambda}^{(h)}B),
 \label{eq:graph-evolution}\\
 \norm{L^p\partial_\lambda^a\partial_v^h\partial_\tau^cS_{\lambda,\tau}}
 &\le\ee^{-4\tau}\norm{L^{p+c}(f_{a,\lambda}^{(h)}B)}.
 \label{eq:graph-evolution-bound}
\end{align}
Strong continuity in the graph spaces, with one extra finite power for
each differentiation, gives the actual right derivatives at zero.
Compactness of the parameterized initial family makes convergence
uniform on $|\lambda|\le M$.

\subsection{Weighted spatial estimates}

Graph control alone does not justify an entire pointwise Fourier
readout. We therefore propagate exponential spatial moments explicitly.
Let
\begin{equation}
 K_\tau(y)=(4\pi\tau)^{-1/2}\ee^{-y^2/(4\tau)},\qquad \tau>0,
 \qquad K_0*f=f.
 \label{eq:heat-kernel}
\end{equation}
The commuting radial and angular multipliers give
\begin{equation}
 S_{\lambda,\tau}(v)
 =(K_\tau*F_\lambda)(v)\,\ee^{-4\tau A}B.
 \label{eq:heat-factorization}
\end{equation}
This is initially an equality in $\HH$. The following estimates also
place both sides in all weighted spaces needed for observation.

The weight $w_{r,b}(v)=(1+|v|)^r\ee^{b|v|}$ is submultiplicative:
$w_{r,b}(x+y)\le w_{r,b}(x)w_{r,b}(y)$. For a standard real Gaussian
$X$, set
\begin{equation}
 \gamma_{r,b}(T)=
 \mathbb E\bigl[(1+\sqrt{2T}|X|)^r\ee^{b\sqrt{2T}|X|}\bigr],
 \qquad \gamma_{r,b}(0)=1.
 \label{eq:gamma}
\end{equation}
Then $\norm{w_{r,b}K_\tau}_1\le\gamma_{r,b}(T)<\infty$ for
$0<\tau\le T$. For an explicit bound, write
$m_\ell=\mathbb E|X|^\ell
=2^{\ell/2}\Gamma((\ell+1)/2)/\sqrt\pi$. Exponential tilting of
the Gaussian gives
\begin{equation}
 \gamma_{r,b}(T)\le2\ee^{b^2T}
 \sum_{j=0}^r\binom rj\sum_{\ell=0}^j\binom j\ell
 (2bT)^{j-\ell}(2T)^{\ell/2}m_\ell.
 \label{eq:gamma-bound}
\end{equation}
To see this, expand $(1+|Y|)^r$ for $Y=\sqrt{2T}X$, bound
$\ee^{b|Y|}$ by $\ee^{bY}+\ee^{-bY}$, and complete the square.
The tilted means are $\pm2bT$, so the triangle and binomial inequalities
give \eqref{eq:gamma-bound}. The convention $0^0=1$ covers the endpoints;
the bound need not be sharp there.

Tonelli's theorem and submultiplicativity now imply
\begin{equation}
 \norm{w_{r,b}(K_\tau*f)}_1
 \le\gamma_{r,b}(T)\norm{w_{r,b}f}_1,
 \qquad 0\le\tau\le T.
 \label{eq:weighted-convolution}
\end{equation}
This estimate propagates every fixed polynomial and exponential
moment. No claim that positive-time convolution preserves the source's
original double-exponential spatial decay is needed.

Move derivatives of the heat kernel onto the smooth, integrable prepared
data. For any finite $a,c,h\ge0$, this yields
\begin{equation}
 \partial_\lambda^a\partial_\tau^c\partial_v^hS_{\lambda,\tau}
 =\sum_{j=0}^c\binom cj(-4)^{c-j}
 (K_\tau*f_{a,\lambda}^{(h+2j)})\,
 A^{c-j}\ee^{-4\tau A}B.
 \label{eq:heat-derivatives}
\end{equation}
Consequently, on the parameter window \eqref{eq:windows},
\begin{align}
 &\norm{w_{r,b}\partial_\lambda^a\partial_\tau^c
                  \partial_v^hS_{\lambda,\tau}}_{L^1(\HS)}
 \le\mathcal M_{a,c,h;r,b}(M,T),
 \label{eq:heat-majorant}\\
 &\mathcal M_{a,c,h;r,b}(M,T)
 :=\gamma_{r,b}(T)\sum_{j=0}^c\binom cj4^{c-j}
       \sqrt{\mu_{2(c-j)}}\,
       \mathcal N_{a,h+2j;r,b}(M)<\infty.
 \label{eq:M-definition}
\end{align}
An additional left power $L^p$ has the same absolute estimate with $c$
replaced by $c+p$. An angular tail replaces the appropriate $\mu$
moment by the vanishing tail in \eqref{eq:bank-tail}. These estimates
justify complete angular and reference limits as well as the radial
integrals.

\subsection{Why the limits at zero are actual derivatives}

For $\tau>0$, integration by parts and the Gaussian kernel justify
\eqref{eq:heat-derivatives}. A bound alone would not establish that
its limit at zero differentiates the original field. The missing step
is continuity in the weighted Banach spaces.

Translation is continuous in $L^1(w_{r,b}\dd v)$. Indeed, approximate
by compactly supported continuous functions; for the remainder, the
translation norm is at most $w_{r,b}(y)$ times its original norm.
Writing $K_\tau*f-f$ as the Gaussian average of translation differences
therefore gives convergence to zero. A dominating random variable is
\begin{equation*}
 \bigl(1+w_{r,b}(\sqrt{2T}X)\bigr)\norm{w_{r,b}f}_1,
\end{equation*}
which is integrable by \eqref{eq:gamma}. This convergence is uniform on
compact sets of $f$ in the weighted space. Section~\ref{sec:source}
established precisely that compactness for
$\{f_{a,\lambda}^{(h)}:|\lambda|\le M\}$.

On the angular side,
$A^s\ee^{-4\tau A}B\to A^sB$ in Hilbert--Schmidt norm by
\eqref{eq:bank-moments} and dominated convergence. Hence every
expression on the right side of \eqref{eq:heat-derivatives} has a
continuous weighted limit at zero. Integrating the positive-time
identity
\begin{equation*}
 \partial_\tau(K_\tau*f)=K_\tau*f''
\end{equation*}
from a positive lower endpoint and then taking that endpoint to zero
shows that its limiting value is the actual right derivative. Iteration
on the controlled derivatives gives every finite order. The same
argument, using the next parameter insertion or spatial derivative,
proves all stated mixed interchanges. Thus
\eqref{eq:heat-majorant} controls continuous integral identities,
including the boundary $\tau=0$.

\section{The calibrated response and its derivatives}
\label{sec:readout}

\subsection{A well-defined phase-sensitive observation}

For $X\in\HS(\Href,\K)$, the product $B^*X$ is trace class and
\begin{equation}
 \norm{B^*X}_1\le\norm B_{\HS}\norm X_{\HS}=\norm X_{\HS},
 \qquad |\Tr_{\Href}(B^*X)|\le\norm X_{\HS}.
 \label{eq:trace-pairing}
\end{equation}
For example, expanding in orthonormal bases and applying Cauchy--Schwarz
to the two Hilbert--Schmidt coefficient arrays gives the trace bound;
the trace-norm statement is the corresponding product inequality.
This is the Hilbert--Schmidt pairing, conjugate-linear in the fixed
first slot. It is independent of the choice of reference basis. The
matched unitary $J$ fixes the relative preparation and readout phase.

The weighted estimates give a Bochner integral
\begin{equation}
 \widehat S_{\lambda,\tau}(z)
  =(2\pi)^{-1/2}\int_\R\ee^{-\ii zv}S_{\lambda,\tau}(v)\dd v,
 \qquad z\in\C,
 \label{eq:field-fourier}
\end{equation}
and make it entire. For real $k$ this is a continuous representative
of the $L^2$ Fourier transform. Fixed-$k$ evaluation is not a bounded
functional on arbitrary radial $L^2$ data; the weighted regularity of
the prepared family is what makes \eqref{eq:field-fourier} available.

The dimensionless calibrated observation is defined by
\begin{equation}
 \Obs_{\lambda,\tau}(z)
 =\sqrt{\frac\pi2}\,
   \Tr_{\Href}\bigl(B^*\widehat S_{\lambda,\tau}(z)\bigr).
 \label{eq:readout}
\end{equation}
The normalization is absolute. In the supplied physical calibration,
let $S^{\mathrm{phys}}=gS$ with a known $g>0$, and let the raw
phase-sensitive detector output be
\begin{equation}
 y(k)=\eta\Tr_{\Href}(B^*\widehat S^{\mathrm{phys}}(k)),
 \qquad \eta\in\C\setminus\{0\}
 \quad\hbox{known, including its phase}.
 \label{eq:raw-readout}
\end{equation}
Then $\sqrt{\pi/2}\,y/(\eta g)$ is exactly \eqref{eq:readout}.
This law is an explicit model assumption. It is a linear amplitude
observation and is not derived here from a Born probability or a
laboratory detector construction.

There is no residual metric factor. Under the product-to-metric unitary
$U=2$, the prepared field and reference become $US$ and $UB$.
The two factors of $2$ in their pairing cancel the angular metric
measure $1/4$, so the pairing in \eqref{eq:readout} is unchanged.
Nor does the adjoint $B^*$ introduce right-hand heat evolution.

\subsection{Evaluation on the complete spectrum}

\begin{theorem}[Exact response]
\label{thm:response}
For the preparation \eqref{eq:prepared-field}, left evolution
\eqref{eq:evolved-field}, and calibrated readout \eqref{eq:readout},
\begin{equation}
 \Obs_{\lambda,\tau}(z)
  =q(\tau)\ee^{-\tau z^2}H_\lambda(z),
 \qquad q(\tau)=\frac{Q_1(\pi+4\tau)}{Q_1(\pi)}>0.
 \label{eq:response}
\end{equation}
This identity has all the finite-order regularity and uniform compact
window bounds stated in Theorem~\ref{thm:overview}.
\end{theorem}

\begin{proof}
For real $k$, the complete multiplier gives
\begin{equation}
 P_n\widehat S_{\lambda,\tau}(k)
 =\ee^{-\tau k^2}\ee^{-4\tau n^2}
     \widehat F_\lambda(k)P_nB.
 \label{eq:channel-response}
\end{equation}
The trace is absolutely convergent by \eqref{eq:trace-pairing} and the
bank moments. It can therefore be evaluated either shell by shell or
by cyclicity of the trace-class products:
\begin{align}
 \Tr_{\Href}(B^*\ee^{-4\tau A}B)
 &=\Tr_\K(\ee^{-4\tau A}R)
 \nonumber\\
 &=\frac1Q\sum_{n\ge1}\ee^{-(\pi+4\tau)n^2}
  =q(\tau).
 \label{eq:trace-q}
\end{align}
There is one left heat factor, which is why the shift is $4\tau$.
Combining \eqref{eq:channel-response} with
\eqref{eq:H-fourier} and \eqref{eq:readout} cancels
$\sqrt{\pi/2}\sqrt{2/\pi}=1$ and proves \eqref{eq:response} on
the real axis. Both sides are entire in $z$, so the identity theorem
extends it to $\C$. The positivity of $q$ follows termwise from its
complete series. The derivative claims follow from the bounds below.
\end{proof}

At $\tau=0$, $q(0)=1$ and $\Obs_{\lambda,0}=H_\lambda$.
In particular $\Obs_{0,0}=\xi((1+\ii z)/2)/8$. This checks the
source, Fourier, bank, volume, and detector normalizations together.
The source sum and the physical-shell sum enter different parts of
the calculation; no diagonal identification between them is made.

\subsection{Integral and scalar forms of every finite derivative}

Combining the two normalization factors directly in
\eqref{eq:readout} gives the useful integral formula
\begin{equation}
 \partial_\lambda^a\partial_\tau^c\partial_z^d\Obs_{\lambda,\tau}(z)
 =\frac12\Tr_{\Href}\left[
 B^*\int_\R(-\ii v)^d\ee^{-\ii zv}
       \partial_\lambda^a\partial_\tau^cS_{\lambda,\tau}(v)\dd v
 \right].
 \label{eq:readout-derivative}
\end{equation}
For $|\Im z|\le b$, its absolute value is at most
\begin{equation}
 \frac12\mathcal M_{a,c,0;d,b}(M,T),
 \label{eq:readout-majorant}
\end{equation}
because $|v|^d\ee^{b|v|}\le w_{d,b}(v)$. The continuity results of
Section~\ref{sec:heat} justify every interchange and the actual
right-time derivatives at zero. This also proves the entire extension
used in the theorem without relying on a formal complex Fourier
evaluation on generic $L^2$ vectors.

There is a complementary scalar recursion for calculating these
derivatives. The angular series gives
\begin{equation}
 q^{(s)}(\tau)=\frac{(-4)^s}{Q}
       \sum_{n\ge1}n^{2s}\ee^{-(\pi+4\tau)n^2},
 \qquad |q^{(s)}(\tau)|\le4^s\mu_s.
 \label{eq:q-derivatives}
\end{equation}
Define polynomials in $(\tau,z)$ by
\begin{equation}
 U_{0,r}(\tau,z)=(-z^2)^r,\qquad
 U_{j+1,r}=\partial_zU_{j,r}-2\tau zU_{j,r}.
 \label{eq:U-recursion}
\end{equation}
Then
$\partial_z^j\partial_\tau^r\ee^{-\tau z^2}
=\ee^{-\tau z^2}U_{j,r}(\tau,z)$.
Writing $h_a=\partial_\lambda^aH_\lambda$, finite Leibniz expansion
of \eqref{eq:response} yields
\begin{equation}
 \partial_\lambda^a\partial_\tau^c\partial_z^d\Obs
 =\ee^{-\tau z^2}\sum_{s=0}^c\sum_{j=0}^d
   \binom cs\binom dj q^{(s)}(\tau)
   U_{j,c-s}(\tau,z)\partial_z^{d-j}h_a(z).
 \label{eq:scalar-mixed}
\end{equation}
Let $U^\#_{j,r}(T,\rho)$ be the polynomial obtained by replacing all
coefficients of $U_{j,r}$ by their absolute values and then evaluating
at $T,\rho$. On \eqref{eq:windows}, set $\rho=\sqrt{Z^2+b^2}$.
Since $|\ee^{-\tau z^2}|\le\ee^{Tb^2}$, a second explicit bound is
\begin{equation}
 |\partial_\lambda^a\partial_\tau^c\partial_z^d\Obs|
 \le\frac{B_0\ee^{Tb^2}}2
 \sum_{s=0}^c\sum_{j=0}^d\binom cs\binom dj
 4^s\mu_s U^\#_{j,c-s}(T,\rho)J_{2a+d-j,2}(M,b).
 \label{eq:scalar-mixed-bound}
\end{equation}
Every sum here is finite. Neither this formula nor
\eqref{eq:readout-majorant} asserts a bound uniform in derivative order,
an infinite derivative expansion, or evolution at negative or complex
proper time.

\section{Zeros, conjugated dynamics, and the scope of recovery}
\label{sec:consequences}

\subsection{Exact zero preservation}

For fixed $\tau\ge0$, define
\begin{equation}
 m_\tau(z)=q(\tau)\ee^{-\tau z^2}.
 \label{eq:multiplier}
\end{equation}
It is entire and nowhere zero. Hence if, for some fixed $\lambda$,
$H_\lambda(z)=(z-z_0)^r h(z)$ with $h(z_0)\ne0$, then
\begin{equation*}
 \Obs_{\lambda,\tau}(z)=(z-z_0)^r m_\tau(z)h(z)
\end{equation*}
has exactly the same multiplicity $r$ at $z_0$. Since $H_\lambda(0)>0$,
the identically zero case does not occur. Thus zeros, multiplicities,
and their collision loci as $\lambda$ varies are preserved exactly.
This conclusion follows from invertible multiplication; it supplies
no new inequality constraining where those collisions can occur.

\subsection{The conjugated Newman equation}

Introduce the differential operator
\begin{equation}
 \mathscr D_\tau=\partial_z+2\tau z.
 \label{eq:covariant-derivative}
\end{equation}
For any differentiable $h$,
\begin{equation}
 \mathscr D_\tau(m_\tau h)=m_\tau h'.
 \label{eq:conjugation}
\end{equation}
Iterating this identity and using \eqref{eq:newman-heat} gives
\begin{equation}
 \boxed{\partial_\lambda\Obs_{\lambda,\tau}
       =-\mathscr D_\tau^2\Obs_{\lambda,\tau}}
 \label{eq:conjugated-heat}
\end{equation}
or, in ordinary derivatives,
\begin{equation}
 \partial_\lambda\Obs
 =-\partial_z^2\Obs-4\tau z\partial_z\Obs
                 -(2\tau+4\tau^2z^2)\Obs.
 \label{eq:conjugated-heat-expanded}
\end{equation}
The term $2\tau\Obs$ comes from differentiating the coefficient
$2\tau z$ inside $\mathscr D_\tau^2$ and must be retained.
For every finite $a$,
$\partial_\lambda^a\Obs=(-1)^a\mathscr D_\tau^{2a}\Obs$;
these are finite differential identities justified by the preceding
regularity, not claims about infinite series of operators.

The proper-time equation is instead
\begin{equation}
 \partial_\tau\Obs_{\lambda,\tau}(z)
 =\left(\frac{q'(\tau)}{q(\tau)}-z^2\right)
                         \Obs_{\lambda,\tau}(z).
 \label{eq:proper-time-response}
\end{equation}
Equations~\eqref{eq:conjugated-heat} and
\eqref{eq:proper-time-response} show directly why Newman deformation
and auxiliary proper-time evolution are different operations.

\subsection{Inversion and the first derivative pair}

The exact inverse is
\begin{equation}
 H_\lambda(z)=q(\tau)^{-1}\ee^{\tau z^2}
                         \Obs_{\lambda,\tau}(z).
 \label{eq:inverse}
\end{equation}
For real $k$ and $\tau>0$, its absolute multiplier is
$q(\tau)^{-1}\ee^{\tau k^2}$, which tends to infinity as
$|k|\to\infty$. Hence the construction does not give an inverse
bounded uniformly over all real frequencies. On a fixed window
$|k|\le K$, $0\le\tau\le T$, the scalar inverse is bounded by
$q(T)^{-1}\ee^{TK^2}$ because $q$ is positive and decreasing.
This is a finite-window attenuation calculation, not a source-specific
collision estimate.

The pair of an amplitude and its first derivative is the natural next
diagnostic. Fix $s>0$ to specify the relative units of derivative error
and value error, and define, at a real frequency,
\begin{equation}
 X=\begin{pmatrix}H_\lambda(k)\\sH_\lambda'(k)\end{pmatrix},
 \qquad
 Y=\begin{pmatrix}\Obs_{\lambda,\tau}(k)\\
                         s\partial_k\Obs_{\lambda,\tau}(k)\end{pmatrix}.
\end{equation}
Differentiating \eqref{eq:response} once gives the exact algebraic map
\begin{equation}
 Y=mP_aX,\qquad m=q(\tau)\ee^{-\tau k^2},\qquad
 a=2s\tau k,\qquad
 P_a=\begin{pmatrix}1&0\\-a&1\end{pmatrix}.
 \label{eq:first-jet}
\end{equation}
Equivalently,
\begin{equation}
 P_a^{-1}Y=
 \begin{pmatrix}\Obs\\s(\partial_k+2\tau k)\Obs\end{pmatrix}
 =mX.
 \label{eq:covariant-pair}
\end{equation}
This diagonal form does not remove measurement error by itself. Under
the same change of coordinates an additive error set $E$ becomes
$P_a^{-1}E$. Similarly, source uncertainty $E$ propagated through the
exact map $\mathsf A=mP_a$ and then inverted returns
$\mathsf A^{-1}(\mathsf A E)=E$.
Changing the error model after changing coordinates would compare
different information. These elementary observations motivate the
pending sharp conditioning task in Section~\ref{sec:outstanding}; no
sharp first-jet recovery theorem is claimed here.

\subsection{What the representation establishes}

The construction proves a complete analytic realization of an exact
response under specified preparation and observation laws. It retains
the geometric spectrum, every physical angular shell, the full radial
integral, the passive reference, and the absolute gain. Its validity
does not rest on numerical truncation. It does not derive the source
scan or imprint from free cylinder dynamics, nor does it derive the
linear phase-sensitive detector from an intrinsic measurement model.

The positivity of $L$, $R$, and $\Phi$ also has distinct meanings.
The Hamiltonian is bounded below, the bank is a positive normalized
angular operator, and the source is positive. A real-frequency
observation nevertheless integrates the source against an oscillatory
cosine. Its zeros describe cancellation in this amplitude, not a zero
eigenvalue of the positive Hamiltonian. Hilbert-space positivity alone
therefore does not exclude a multiple zero of $H_\lambda$.

The existing control families make this distinction concrete
\cite{CollisionControls}. Their role is to test the premises of a
proposed estimate, rather than to replace the canonical source:
\begin{itemize}
\item For applicable Euler-factor controls, smooth positive-source
Fourier and heat operations persist, while the unchanged canonical
theta arithmetic fails. A generic representation of this kind rejects
none of their recorded collisions.
\item Weighted-square and finite-prime constructions can retain
positive half-line formulas while losing complete unit coefficients
and the canonical reflection identity. Matching finitely many endpoint
conditions does not establish the global smooth domains used here.
\item Applicable reflected super-Gaussian or smooth positive primitive
controls retain smoothness, positivity, reflection, and exponential
moments without the exact theta expansion. Where the analytic
hypotheses hold, the same nonvanishing multiplier preserves their
collisions as well.
\item Atomic or finite-profile controls can have Fourier observations
and moments without satisfying the smooth field-domain hypotheses.
No smooth-field realization for them follows from the present proof.
\end{itemize}
Consequently, naming the canonical source identifies the input but
does not yet quantitatively separate it from colliding controls.
Any future collision estimate must identify the additional arithmetic
or analytic property that makes such a separation possible.

\section{Outstanding tasks}
\label{sec:outstanding}

The established endpoint is the complete calibrated response theorem,
together with its domains, finite-order estimates, exact zero
preservation, and conjugated dynamics. The next work concerns what can
be recovered from that response and whether the canonical source
supports a further estimate.

\paragraph{Sharp finite-window first-jet conditioning.}
The decision recorded as TND-A-B02 \cite{TNDAuxConditioning} and the
frozen TND-A-E03 request \cite{TNDAuxE03Request} ask for the exact
worst-case recovery error for \eqref{eq:first-jet}. The parameters
$M,K,T$ are finite, $s>0$ is fixed, and the norms and error units must
be stated before comparing proper times. For an additive Euclidean
error ball on the observed pair, the question is to determine the
sharp pointwise and window bounds governed by
\begin{equation}
 \norm{(mP_a)^{-1}}_2
 =q(\tau)^{-1}\ee^{\tau k^2}\norm{P_a^{-1}}_2,
 \qquad |k|\le K,\quad 0\le\tau\le T.
 \label{eq:pending-conditioning}
\end{equation}
Sharpness over unrestricted pair errors must be distinguished from
sharpness over perturbations realizable by admissible theta sources
or globally consistent differentiated error functions. The same
analysis must track the transformed error ellipsoid in
\eqref{eq:covariant-pair} and compare it with uncertainty propagated
from the same source. The present manuscript supplies the exact map;
the requested sharp theorem has not been executed or assessed here.

\paragraph{Assessment of the conditioning result.}
The planned TND-A-E04 assessment should check the completed E03 result
against its fixed error model, endpoints, and sharpness class. Neither
a convenient coordinate choice nor a smaller plotted amplitude is
evidence of improved recovery. This diagnostic has no assumed
physical noise law and no built-in positive collision margin.

\paragraph{A source-specific consequence.}
A consequential next estimate would have to use the exact canonical
theta structure quantitatively. In particular, the present positivity,
regularity, reflection, and heat identities do not by themselves bound
the simultaneous smallness of $H_\lambda$ and $H_\lambda'$ near a
potential collision. A proposed lower bound or signed inequality must
state its parameter range and identify precisely which premise fails
for the relevant colliding controls. If the conditioning diagnostic
shows only attenuation and reversal, that would justify stopping this
representation as an estimating route unless an additional
source-specific mechanism is found.

\paragraph{Intrinsic preparation and physical implementation.}
The scan $\beta(u)=\pi\ee^{4u}$, processing operator
$(\partial_u^2-1)/8$, spatial imprint, linear amplitude gain, matched
passive bank, and phase-sensitive detector law are supplied. Deriving
any of them from a specified physical apparatus would be a separate
task with additional assumptions. Such a derivation should also
distinguish Euclidean evolution from real unitary laboratory time and
account for preparation, calibration, and detection errors.

\paragraph{Independent mathematical review and supporting evidence.}
The current proof notes and their assessments are internal LLM-assisted
work with shared premises. Independent specialist review should
examine the geometric conventions, complete-domain construction,
weighted zero-time argument, and absolute calibration before stronger
reliance is placed on the representation. Numerical checks may support
finite-window constants and truncation bounds after the analytic
questions are fixed. They would remain supporting evidence: a finite
computation cannot replace a uniform estimate or a collision exclusion.


\section*{Acknowledgements and provenance}
This manuscript was drafted for Edward Baker with substantial assistance
from OpenAI's \texttt{gpt-6-astra} through Codex, with recorded reasoning
effort \texttt{max}. LLM assistance included mathematical exposition,
comparison with the proof notes and primary sources, local algebraic and
analytic checks, and document production. The principal proof notes and
their internal assessments were prepared with \texttt{gpt-6.1-sol}, effort
\texttt{ultra}. Those assessments have shared premises and include work
in the proof-author chat; they are not independent specialist review or
formal verification. The present synthesis preserves that evidence status.
No numerical experiment is used as a premise of any result.

\begingroup
\small
\bibliographystyle{plainurl}
\bibliography{hamiltonian_auxiliary_sources}
\endgroup
\end{document}
